%  LaTeX support: latex@mdpi.com 
%  For support, please attach all files needed for compiling as well as the log file, and specify your operating system, LaTeX version, and LaTeX editor.

%=================================================================
\documentclass[jmse,article,submit,pdftex,moreauthors]{Definitions/mdpi} 

%=================================================================
% MDPI internal commands
\firstpage{1} 
\makeatletter 
\setcounter{page}{\@firstpage} 
\makeatother
\pubvolume{1}
\issuenum{1}
\articlenumber{0}
\pubyear{2023}
\copyrightyear{2023}
%\externaleditor{Academic Editor: Firstname Lastname}
\datereceived{15 March 2023} 
\daterevised{} 
\dateaccepted{} 
\datepublished{} 
\hreflink{https://doi.org/} % If needed use \linebreak
%\doinum{}

%=================================================================
% Add packages and commands here. The following packages are loaded in our class file: fontenc, inputenc, calc, indentfirst, fancyhdr, graphicx, epstopdf, lastpage, ifthen, lineno, float, amsmath, setspace, enumitem, mathpazo, booktabs, titlesec, etoolbox, tabto, xcolor, soul, multirow, microtype, tikz, totcount, changepage, attrib, upgreek, cleveref, amsthm, hyphenat, natbib, hyperref, footmisc, url, geometry, newfloat, caption

\graphicspath{{figures/}} % path to folder with figures
\usepackage{subcaption}
\usepackage{xcolor}
\usepackage{gensymb} % for \textdegree
\usepackage{tabularx}

\usepackage{listings}
\usepackage{xcolor}
\definecolor{deepblue}{rgb}{0,0,0.5}
\definecolor{deepred}{rgb}{0.6,0,0}
\definecolor{deepmagenta}{RGB}{195,12,214}
\definecolor{deepgreen}{RGB}{29,122,30}
\definecolor{mintgreen}{RGB}{192,249,192}
\definecolor{codepurple}{RGB}{204, 0, 153}
\definecolor{LightCyan}{rgb}{0.88,1,1}
\definecolor{GainsboroPurple}{RGB}{247,176,247}
\definecolor{LightSlateGrey}{RGB}{124,118,118}
%    
\lstdefinestyle{mystyle}{
    commentstyle=\color{LightSlateGrey},
    keywordstyle=\color{deepgreen}\bfseries,
    emphstyle=\color{deepred},
    numberstyle=\tiny\color{deepmagenta},
    stringstyle=\color{codepurple},
    basicstyle=\ttfamily\footnotesize,
    breakatwhitespace=false,         
    breaklines=true,                 
    captionpos=b,                    
    keepspaces=true,                 
    numbers=left,                    
    numbersep=5pt,                  
    showspaces=false,                
    showstringspaces=false,
    showtabs=false,                  
    tabsize=2
}
\lstset{style=mystyle}

%=================================================================
% Full title of the paper (Capitalized)
\Title{Using GRASS GIS for Computing Vegetation Indices from the Satellite Images for Environmental Monitoring of Mangrove Forests in the Coastal Landscapes of Niger Delta, Nigeria}

% MDPI internal command: Title for citation in the left column along the Southwestern Coast of Nigeria
\TitleCitation{Using GRASS GIS for Computing Vegetation Indices from the Satellite Images for Environmental Monitoring of Mangrove Forests in the Coastal Landscapes of Niger Delta, Nigeria}

% Author Orchid ID: enter ID or remove command
\newcommand{\orcidauthorA}{0000-0002-5759-1089} % Polina Add \orcidA{} behind the author's name
\newcommand{\orcidauthorB}{0000-0002-6461-1551} % Olivier Add \orcidB{} behind the author's name

% Authors, for the paper (add full first names)
\Author{Polina Lemenkova $^{1}$*\orcidA{} and Olivier Debeir $^{1}$\orcidB{}}

%\longauthorlist{yes}

% MDPI internal command: Authors, for metadata in PDF
\AuthorNames{Polina Lemenkova and Olivier Debeir}

% MDPI internal command: Authors, for citation in the left column
\AuthorCitation{Lemenkova, P.; Debeir, O.}
% If this is a Chicago style journal: Lastname, Firstname, Firstname Lastname, and Firstname Lastname.

% Affiliations / Addresses (Add [1] after \address if there is only one affiliation.)
\address{%
$^{1}$ \quad Laboratory of Image Synthesis and Analysis (LISA), École polytechnique de Bruxelles (Brussels Faculty of Engineering), Université Libre de Bruxelles (ULB). Building L, Campus du Solbosch, ULB -- LISA CP165/57, Avenue Franklin D. Roosevelt 50, Brussels 1050, Belgium.
}

% Contact information of the corresponding author
\corres{Correspondence: polina.lemenkova@ulb.be; Tel.: +32 471 86 04 59 (P.L.)}

% Current address and/or shared authorship
%\firstnote{Current address: Affiliation 3.} 

% Abstract (Do not insert blank lines, i.e. \\) 
\abstract{Besides, changes in urban spaces of the major cities caused by the urbanisation in this region (Onitsha, Sapele, Warri, Benin CIty and smaller villages in the Delta State) were evaluated and visualised. }

% Keywords
\keyword{Image processing; Nigeria; West Africa; remote sensing; cartography; Niger River; Niger Delta; R language; programming; climate change} 

% The fields PACS, MSC, and JEL may be left empty or commented out if not applicable
\PACS{91.10.Da; 91.10.Jf; 91.10.Sp; 91.10.Xa; 96.25.Vt; 91.10.Fc; 95.40.+s; 95.75.Qr; 95.75.Rs; 42.68.Wt}
\MSC{86A30; 86-08; 86A99; 86A04}
\JEL{Y91; Q20; Q24; Q23; Q3; Q01; R11; O44; O13; Q5; Q51; Q55; N57; C6; C61}

%%%%%%%%%%%%%%%%%%%%%%%%%%%%%%%%%%%%%%%%%%
\begin{document}

\section{Introduction}

\begin{figure}[H]
\centering
	\includegraphics[width=14.0 cm]{Fig_01.jpg}
	\caption{Topographic map of Nigeria. Software: Generic Mapping Tools (GMT) version 6.1.1, \cite{Wessel}. Data source: GEBCO/SRTM \cite{GEBCO,SRTM}. Cartography source: authors.
	\label{fig01}}
\end{figure}

%%%%%%%%%%%%%%%%%%%%%%%%%%%%%%%%%%%%%%%%%%
\section{Materials and Methods}

The image processing has been performed using the GRASS GIS, version 8.2.1, \href{https://grass.osgeo.org/}{https://grass.osgeo.org/}.
%----------- subsection ----------->
\pagebreak
\subsection{Data}

\begin{figure}[H]
\makebox[0.90\linewidth][r]{%
	\begin{subfigure}[b]{.40\textwidth}
		\centering
			\includegraphics[width=0.87\textwidth]{Fig_03a.jpg}
		\caption{2013}
	\end{subfigure}%
	\begin{subfigure}[b]{.40\textwidth}
		\centering
		\includegraphics[width=0.87\textwidth]{Fig_03b.jpg}
		\caption{2015}
	\end{subfigure}%
	\begin{subfigure}[b]{.40\textwidth}
		\centering
		\includegraphics[width=0.87\textwidth]{Fig_03c.jpg}
		\caption{2018}
	\end{subfigure}%
}\\
\vfill \vspace{1mm}
\makebox[0.90\linewidth][r]{%
	\begin{subfigure}[b]{.40\textwidth}
		\centering
			\includegraphics[width=0.87\textwidth]{Fig_03d.jpg}
		\caption{2020}
	\end{subfigure}%
	\begin{subfigure}[b]{.40\textwidth}
		\centering
		\includegraphics[width=0.87\textwidth]{Fig_03e.jpg}
		\caption{2021}
	\end{subfigure}%
	\begin{subfigure}[b]{.40\textwidth}
		\centering
		\includegraphics[width=0.87\textwidth]{Fig_03f.jpg}
		\caption{2022}
	\end{subfigure}%
}\caption{Landsat 8-9 OLI/TIRS images of Niger Delta in natural colour RGB values showing mangrove forests of the  for six years: (\textbf{a}) 2013/12/20, (\textbf{b}) 2015/12/26, (\textbf{c}) 2018/02/21, (\textbf{d}) 2020/01/22, (\textbf{e}) 2021/12/18, (\textbf{f}) 2022/12/29.}\label{fig:03}
\end{figure}

%----------- subsection ----------->

%----------- subsection ----------->
\subsection{Data preprocessing}

GRASS GIS code used for data preprocessing is presented in Listing \ref{lst01}. It shows the steps of importing Landsat channels by gdal and auxiliary modules and converting the DN pixel values to reflectance values in order to get rid of the haze effect on the original scenes by correction for the top-of-atmosphere radiance of the bands using the \emph{i.landsat.toar} module. The NDVI was computed using the processed bands by the \emph{i.vi} module. 

\begin{lstlisting}[language=bash,caption=GRASS GIS script - 1 for converting the DN pixel values to reflectance for correction of the top-of-atmosphere radiance,style=mystyle,label={lst01}]
#!/bin/sh
# raster metadata:
r.info -r LC08_L2SP_189056_20131220_20200912_02_T1_SR_B1
# checking the selected random band
gdalinfo LC08_L2SP_189056_20131220_20200912_02_T1_SR_B7.TIF
# importing the image subset with 7 Landsat bands and display the raster map
r.in.gdal LC08_L2SP_189056_20131220_20200912_02_T1_SR_B1.TIF out=L8_2013_01
r.in.gdal LC08_L2SP_189056_20131220_20200912_02_T1_SR_B2.TIF out=L8_2013_02
r.in.gdal LC08_L2SP_189056_20131220_20200912_02_T1_SR_B3.TIF out=L8_2013_03
r.in.gdal LC08_L2SP_189056_20131220_20200912_02_T1_SR_B4.TIF out=L8_2013_04
r.in.gdal LC08_L2SP_189056_20131220_20200912_02_T1_SR_B5.TIF out=L8_2013_05
r.in.gdal LC08_L2SP_189056_20131220_20200912_02_T1_SR_B6.TIF out=L8_2013_06
r.in.gdal LC08_L2SP_189056_20131220_20200912_02_T1_SR_B7.TIF out=L8_2013_07
r.in.gdal LC08_L2SP_189056_20131220_20200912_02_T1_ST_B8.TIF out=L8_2013_08
r.in.gdal LC08_L2SP_189056_20131220_20200912_02_T1_ST_EMSD.TIF out=L8_2013_09 # Cirrus
r.in.gdal LC08_L2SP_189056_20131220_20200912_02_T1_ST_B10.TIF out=L8_2013_10 # TIRS 1
r.in.gdal LC08_L2SP_189056_20131220_20200912_02_T1_ST_TRAD.TIF out=L8_2013_11 # TIRS 2
# listing the files
g.list rast
# preprocessing:
# copying the Landsat bands to match the input structure of the i.landsat.toar
g.copy raster=L8_2013_01,lsat8_2013.1
g.copy raster=L8_2013_02,lsat8_2013.2
g.copy raster=L8_2013_03,lsat8_2013.3
g.copy raster=L8_2013_04,lsat8_2013.4
g.copy raster=L8_2013_05,lsat8_2013.5
g.copy raster=L8_2013_06,lsat8_2013.6
g.copy raster=L8_2013_07,lsat8_2013.7
g.copy raster=L8_2013_08,lsat8_2013.8
g.copy raster=L8_2013_09,lsat8_2013.9
g.copy raster=L8_2013_10,lsat8_2013.10
g.copy raster=L8_2013_11,lsat8_2013.11
# converting the DN pixel values to reflectance values using DOS1
i.landsat.toar input=lsat8_2013. output=lsat8_2013_toar. sensor=oli8 \
    method=dos1 date=2013-12-20 sun_elevation=51.85549756 \
    product_date=2013-12-20 gain=HHHLHLHHL
\end{lstlisting}

%------------------------->
\subsection{Normalized Difference Vegetation Index (NDVI)}

\begin{verbatim}
	NDVI = (NIR - Red) / (NIR + Red)
\end{verbatim}

\begin{lstlisting}[language=bash,caption=GRASS GIS script for computing the NDVI,style=mystyle,label={lst01}]
g.region raster=lsat8_2013_toar.4 -p
i.vi red=lsat8_2013_toar.4 nir=lsat8_2013_toar.5 viname=ndvi \
       output=lsat8_2013.ndvi --overwrite
r.colors lsat8_2013.ndvi color=ndvi
d.mon wx0
g.region raster=lsat8_2013_toar.4 -p
d.rast lsat8_2013.ndvi
d.legend raster=lsat8_2013.ndvi range=-1,1 title="NDVI" title_fontsize=14 font=Helvetica fontsize=12 -t -s -b border_color=white thin=12 label_step=0.1 -d
d.out.file output=Nigeria_NDVI_2013 format=jpg --overwrite\end{lstlisting}

%------------------------->
\subsection{Atmospherically Resistant Vegetation Index (ARVI)}

\begin{verbatim}
	ARVI = (NIR - (2.0*Red - Blue)) / ( NIR + (2.0*Red - Blue))
\end{verbatim}

\begin{lstlisting}[language=bash,caption=GRASS GIS script for computing the ARVI,style=mystyle,label={lst01}]
g.region raster=lsat8_2013_toar.3 -p
i.vi blue=lsat8_2013_toar.2 red=lsat8_2013_toar.4 nir=lsat8_2013_toar.5 \
       viname=arvi output=lsat8_2013.arvi --overwrite
r.colors lsat8_2013.arvi color=bcyr -e
d.mon wx0
d.rast lsat8_2013.arvi
d.legend raster=lsat8_2013.arvi range=-0.7,0.3 title="ARVI" title_fontsize=14 font=Helvetica fontsize=12 -t -s -b border_color=white thin=12 label_step=0.1 -d
d.out.file output=Nigeria_ARVI_2013 format=jpg --overwrite
\end{lstlisting}

%------------------------->
\subsection{Green Atmospherically Resistant Vegetation Index (GARI)}

\begin{verbatim}
	GARI = ( NIR - (Green - (Blue - Red))) / ( NIR + (Green - (Blue - Red)))
\end{verbatim} 

\begin{lstlisting}[language=bash,caption=GRASS GIS script for computing the GARI,style=mystyle,label={lst01}]
# GARI = (nirchan - (2.0*redchan - bluechan)) / ( nirchan + (2.0*redchan - bluechan))
g.region raster=lsat8_2013_toar.3 -p
i.vi blue=lsat8_2013_toar.2 green=lsat8_2013_toar.3 red=lsat8_2013_toar.4 \
     nir=lsat8_2013_toar.5 viname=gari output=lsat8_2013.gari --overwrite
r.colors lsat8_2013.gari color=rainbow -e
d.mon wx1
d.rast lsat8_2013.gari
d.legend raster=lsat8_2013.gari range=-0.5,0.7 title="GARI" title_fontsize=14 font="Helvetica" fontsize=12 -t -b bgcolor=white label_step=0.1 border_color=white thin=8 -d
d.out.file output=Nigeria_GARI_2013 format=jpg --overwrite
\end{lstlisting}

%------------------------->
\subsection{Green Vegetation Index (GVI)}

\begin{verbatim}
	GVI = ( -0.2848 * Blue - 0.2435 * Green - 0.5436 * Red + 0.7243 * NIR + 0.0840 * SWIR1- 0.1800 * SWI2)
\end{verbatim}

\begin{lstlisting}[language=bash,caption=GRASS GIS script for computing the GVI,style=mystyle,label={lst01}]
g.region raster=lsat8_2013_toar.3 -p
i.vi blue=lsat8_2013_toar.2 green=lsat8_2013_toar.3 red=lsat8_2013_toar.4 nir=lsat8_2013_toar.5 band5=lsat8_2013_toar.6 band7=lsat8_2013_toar.7 viname=gvi output=lsat8_2013.gvi --overwrite
r.colors.stddev lsat8_2013.gvi
d.mon wx0
d.rast lsat8_2013.gvi
d.legend raster=lsat8_2013.gvi range=-1,0.7 title="GVI" title_fontsize=14 font="Helvetica" fontsize=12 -t -b bgcolor=white label_step=0.1 border_color=white thin=8 -d
d.out.file output=Nigeria_GVI_2013 format=jpg --overwrite
\end{lstlisting}

%------------------------->
\subsection{Difference Vegetation Index (DVI)}

\begin{verbatim}
	DVI = (NIR - Red)
\end{verbatim}

\begin{lstlisting}[language=bash,caption=GRASS GIS script for computing the DVI,style=mystyle,label={lst01}]
g.region raster=lsat8_2013_toar.1 -p
i.vi blue=lsat8_2013_toar.2 red=lsat8_2013_toar.4 nir=lsat8_2013_toar.5 viname=dvi output=lsat8_2013.dvi --overwrite
r.colors lsat8_2013.dvi color=bgyr -e
d.mon wx0
d.rast lsat8_2013.dvi
d.legend raster=lsat8_2013.dvi range=-0.1,0.3 title="DVI" title_fontsize=14 font="Helvetica" fontsize=12 -t -b bgcolor=white label_step=0.02 border_color=white thin=8 -d
d.out.file output=Nigeria_DVI_2013 format=jpg --overwrite
\end{lstlisting}

%------------------------->
\subsection{Perpendicular Vegetation Index (PVI)}

\begin{verbatim}
	PVI = sin(a)NIR-cos(a)Red
\end{verbatim} 

\begin{lstlisting}[language=bash,caption=GRASS GIS script for computing the PVI,style=mystyle,label={lst01}]
g.region raster=lsat8_2013_toar.1 -p
i.vi red=lsat8_2013_toar.4 nir=lsat8_2013_toar.5 viname=pvi output=lsat8_2013.pvi --overwrite
r.colors lsat8_2013.pvi color=bgyr -e
d.mon wx0
d.rast lsat8_2013.pvi
d.legend raster=lsat8_2013.pvi range=-0.1,0.3 title="PVI" title_fontsize=14 font="Helvetica" fontsize=12 -t -b bgcolor=white label_step=0.02 border_color=white thin=8 -d
d.out.file output=Nigeria_PVI_2013 format=jpg --overwrite
\end{lstlisting}

%%%%%%%%%%%%%%%%%%%%%%%%%%%%%%%%%%%%%%%%%%
\section{Results}

%----------- subsection ----------->

\begin{figure}[H]
\makebox[0.90\linewidth][r]{%
	\begin{subfigure}[b]{.55\textwidth}
		\centering
			\includegraphics[width=7.5cm]{Fig_04a.jpg}
		\caption{2013: NDVI}
	\end{subfigure}%
	\begin{subfigure}[b]{.55\textwidth}
		\centering
			\includegraphics[width=7.5cm]{Fig_04b.jpg}
		\caption{2013: ARVI}
	\end{subfigure}%
}\\
\vfill \vspace{1mm}
\makebox[0.90\linewidth][r]{%
	\begin{subfigure}[b]{.55\textwidth}
		\centering
			\includegraphics[width=7cm]{Fig_04c.jpg}
		\caption{2013: GARI}
	\end{subfigure}
	\begin{subfigure}[b]{.55\textwidth}
		\centering
			\includegraphics[width=7cm]{Fig_04d.jpg}
		\caption{2013: GVI}
	\end{subfigure}%
}\\
\vfill \vspace{1mm}
\makebox[0.90\linewidth][r]{%
	\begin{subfigure}[b]{.55\textwidth}
		\centering
		\includegraphics[width=7cm]{Fig_04e.jpg}
		\caption{2013: DVI}
	\end{subfigure}%
	\begin{subfigure}[b]{.55\textwidth}
		\centering
		\includegraphics[width=7cm]{Fig_04f.jpg}
		\caption{2013: PVI}
	\end{subfigure}%
}\caption{NDVI computed from the Landsat 8-9 OLI/TIRS images of Niger Delta for 2013 years: (\textbf{a}) NDVI, (\textbf{b}) ARVI, (\textbf{c}) GARI, (\textbf{d}) GVI, (\textbf{e}) DVI, (\textbf{f}) PVI.}\label{fig:03}
\end{figure}

%----------- subsection ----------->

\begin{figure}[H]
\makebox[0.90\linewidth][r]{%
	\begin{subfigure}[b]{.40\textwidth}
		\centering
			\includegraphics[width=0.87\textwidth]{Fig_05a.jpg}
		\caption{2013}
	\end{subfigure}%
	\begin{subfigure}[b]{.40\textwidth}
		\centering
		\includegraphics[width=0.87\textwidth]{Fig_03b.jpg}
		\caption{2015}
	\end{subfigure}%
	\begin{subfigure}[b]{.40\textwidth}
		\centering
		\includegraphics[width=0.87\textwidth]{Fig_03c.jpg}
		\caption{2018}
	\end{subfigure}%
}\\
\vfill \vspace{1mm}
\makebox[0.90\linewidth][r]{%
	\begin{subfigure}[b]{.40\textwidth}
		\centering
			\includegraphics[width=0.87\textwidth]{Fig_03d.jpg}
		\caption{2020}
	\end{subfigure}%
	\begin{subfigure}[b]{.40\textwidth}
		\centering
		\includegraphics[width=0.87\textwidth]{Fig_03e.jpg}
		\caption{2021}
	\end{subfigure}%
	\begin{subfigure}[b]{.40\textwidth}
		\centering
		\includegraphics[width=0.87\textwidth]{Fig_03f.jpg}
		\caption{2022}
	\end{subfigure}%
}\caption{ARVI computed from the Landsat 8-9 OLI/TIRS images of Niger Delta for six years: (\textbf{a}) 2013/12/20, (\textbf{b}) 2015/12/26, (\textbf{c}) 2018/02/21, (\textbf{d}) 2020/01/22, (\textbf{e}) 2021/12/18, (\textbf{f}) 2022/12/29.}\label{fig:03}
\end{figure}

%----------- subsection ----------->
\subsection{GARI}

%----------- subsection ----------->
\subsection{GEMI}

%----------- subsection ----------->
\subsection{Subsection}

%----------- subsection ----------->
\subsection{Subsection}

%%%%%%%%%%%%%%%%%%%%%%%%%%%%%%%%%%%%%%%%%%
\section{Discussion}

%%%%%%%%%%%%%%%%%%%%%%%%%%%%%%%%%%%%%%%%%%
\section{Conclusions}

%%%%%%%%%%%%%%%%%%%%%%%%%%%%%%%%%%%%%%%%%%
\authorcontributions{Supervision, conceptualization, methodology, software, resources, funding acquisition, and project administration, O.D.; writing---original draft preparation, methodology, software, data curation, visualization, formal analysis, validation, writing---review and editing, and investigation, P.L. All authors have read and agreed to the published version of the manuscript.}

\funding{The publication was funded by the Editorial Office of JMSE, Multidisciplinary Digital Publishing Institute (MDPI), by providing 100\% discount for the APC of this manuscript. This project was supported by the Federal Public Planning Service Science Policy or Belgian Science Policy Office, Federal Science Policy - BELSPO (B2/202/P2/SEISMOSTORM).}

\institutionalreview{Not applicable.}

\informedconsent{Not applicable.}

\dataavailability{Not applicable} 

\acknowledgments{The authors thank the anonymous reviewers for reading, suggestions and comments that improved an earlier version of this manuscript.}

\conflictsofinterest{The authors declare no conflict of interest.} 

\abbreviations{Abbreviations}{
The following abbreviations are used in this manuscript:\\

\noindent 
\begin{tabular}{@{}ll}
GMT & Generic Mapping Tools\\
GRASS GIS & Geographic Resources Analysis Support System Geographic Information System
\end{tabular}
}

%%%%%%%%%%%%%%%%%%%%%%%%%%%%%%%%%%%%%%%%%%
\begin{adjustwidth}{-\extralength}{0cm}
%\printendnotes[custom] % Un-comment to print a list of endnotes

\bibliography{Nigeria.bib}
\reftitle{References}
\nocite{*}

%=====================================
% References, variant A: external bibliography
%=====================================

%%%%%%%%%%%%%%%%%%%%%%%%%%%%%%%%%%%%%%%%%%
\end{adjustwidth}
\end{document}
